\studySubsection{calc-05-differential-applications-geometry}{第 5 讲 一元函数微分学的应用（一）——几何应用}

\begin{knowledgeBox}
教材页码：第 139 页起。

本讲用于整理切线、法线、函数单调性、极值、最值、凹凸性、拐点、曲率及相关几何应用。新增题目时重点区分一阶导数、二阶导数在几何判定中的作用。
\end{knowledgeBox}

\studySubsection{calc-05-k070}{K070 单调区间判定：从定义域到符号表}

\knowledgeAnchor[MATH1-KN-CALC-0117]{单调区间判定}
\noindent{\footnotesize\textbf{稳定引用：} \knowledgeRef[MATH1-KN-CALC-0117]{单调区间判定}}\par

\begin{knowledgeBox}
\textbf{定位与路线：}本讲解决“函数在哪些区间增、哪些区间减”，并把单调性用于根的唯一性、参数范围与不等式。稳定流程是
\[
\boxed{\text{定义域}\to\text{求导}\to\text{找分界点}\to\text{判导数符号}\to\text{写最大单调区间}}.
\]

\textbf{直接前置（首次完整闭环）：}基本求导公式给出 \(f'(x)\)，拉格朗日中值定理把导数符号传到函数增量；本题需要二者把局部“向上/向下”变成区间上的单调性。使用时先在函数有定义且可导的开区间内求 \(f'\)，再用试点、因式符号或不等式判号。例如 \(f(x)=x^2\) 有 \(f'(x)=2x\)，故在 \((-\infty,0]\) 递减、在 \([0,\infty)\) 递增。这里“临界分界点”包括 \(f'=0\) 的驻点、函数有定义但导数不存在的点，以及定义域断点；最后一种点绝不能跨越。

\textbf{严格核心与推导：}若 \(f\) 在区间 \(I\) 上连续、在内部可导，且内部恒有 \(f'(x)>0\)，则任取 \(x_1<x_2\)，拉格朗日中值定理给出
\[
f(x_2)-f(x_1)=f'(\xi)(x_2-x_1)>0,
\]
故 \(f\) 在 \(I\) 上严格递增；\(f'<0\) 同理推出严格递减。若仅有 \(f'\ge0\)，可先推出不减；当导数只在孤立点为零、且任意子区间上不恒为零时，仍可结合中值定理得到严格递增。

\textbf{方法选择：}“求单调区间”先作完整符号表；“证明方程至多一根”优先考察两边之差的严格单调性；“参数使恒增”把问题转成定义域内 \(f'(x)\ge0\)，再检查边界参数是否仍严格。遇绝对值、分段式或根式，先补不可导点与定义域断点。

\textbf{例 1（标准因式判号）}\quad
\textbf{信号：}多项式求单调区间。\textbf{分析：}导数可因式分解。
\textbf{完整解答：}对 \(f=x^3-3x\)，有 \(f'=3(x-1)(x+1)\)，故 \(f'>0\) 于 \((-\infty,-1)\cup(1,\infty)\)，\(f'<0\) 于 \((-1,1)\)。所以最大递增区间为 \((-\infty,-1]\)、\([1,\infty)\)，最大递减区间为 \([-1,1]\)。
\textbf{条件：}多项式处处连续可导。\textbf{易错：}把两个不相邻递增段合并。\textbf{核验：}取试点 \(-2,0,2\)，导数符号依次为 \(+,-,+\)。

\textbf{例 2（不可导点也要分区）}\quad
\textbf{信号：}\(|x|\) 出现尖点。\textbf{分析：}先列 \(x=0\) 为函数有定义的不可导点。
\textbf{完整解答：} \(f(x)=|x|\) 在 \(x<0\) 时 \(f'=-1\)，在 \(x>0\) 时 \(f'=1\)，故在 \((-\infty,0]\) 严格递减，在 \([0,\infty)\) 严格递增。
\textbf{条件：}两侧分别可导，且函数在 \(0\) 连续。\textbf{易错：}只解 \(f'=0\)，从而漏掉 \(0\)。\textbf{核验：}图像为顶点在原点的 V 形。

\textbf{例 3（参数边界）}\quad
\textbf{信号：}要求含参数函数在 \(x>0\) 的单调性。\textbf{分析：}把参数放进导数不等式。
\textbf{完整解答：}设 \(f(x)=x+a/x\,(x>0)\)，则 \(f'=1-a/x^2\)。若 \(a\le0\)，则 \(f'>0\)，函数严格递增；若 \(a>0\)，则在 \((0,\sqrt a]\) 递减、在 \([\sqrt a,\infty)\) 递增。
\textbf{条件：}定义域仅为 \(x>0\)。\textbf{易错：}把负半轴也列入，或漏 \(a=0\)。\textbf{核验：}当 \(a>0\) 时 \(x=\sqrt a\) 两侧导数由负变正。

\textbf{例 4（根的唯一性迁移）}\quad
\textbf{信号：}证明 \(x+\ln x=0\) 只有一根。\textbf{分析：}令左端为函数并证明严格单调，再证存在。
\textbf{完整解答：}令 \(g(x)=x+\ln x\,(x>0)\)，则 \(g'=1+1/x>0\)，故至多一根；又 \(g(1/2)<0\)、\(g(1)>0\)，由连续性和零点定理至少一根，故恰有一根。
\textbf{条件：}在 \([1/2,1]\) 连续，在 \((0,\infty)\) 可导。\textbf{易错：}只证单调便声称“恰有一根”。\textbf{核验：}存在性与唯一性两条证据相互独立。

\textbf{失分防线：}错误写法“令 \(f'=0\)，所得点之间就是单调区间”漏了不可导点和定义域断点；反例 \(1/x\) 的导数在两侧都负，却不能写成跨越 \(0\) 的单调区间。快速检查：先在草稿顶端写 \(D_f\)，再把“零点、不可导点、断点”逐项打勾。

\textbf{考场表达：}“定义域为……；\(f'(x)=\cdots\)。分界点为……。由导数符号表，\(f\) 在……严格增，在……严格减。”参数题另写边界值单独核验。

\textbf{三级掌握标准：}基础过关：能完整列分界点并作符号表；\(130+\)：能把恒单调转成参数不等式并处理边界；满分能力：能用单调性闭合“存在且唯一”、不等式及跨形式迁移。

\textbf{知识连接：}直接后续为 \knowledgeRef[MATH1-KN-CALC-0118]{K071 局部极值}、\knowledgeRef[MATH1-KN-CALC-0119]{K072 闭区间最值}和 \knowledgeRef[MATH1-KN-CALC-0122]{K075 图形描绘}。
\end{knowledgeBox}

\studySubsection{calc-05-k071}{K071 局部极值：必要条件、变号判据与退化点}

\knowledgeAnchor[MATH1-KN-CALC-0118]{局部极值与必要/充分条件}
\noindent{\footnotesize\textbf{稳定引用：} \knowledgeRef[MATH1-KN-CALC-0118]{局部极值与必要/充分条件}}\par

\begin{knowledgeBox}
\textbf{定位与路线：}局部极值只比较某点附近的函数值。本讲区分“极值点、极值、驻点”，并建立最稳的导数变号判据。

\textbf{直接前置四要素桥：}\knowledgeRef[MATH1-KN-CALC-0117]{K070 单调区间}说明导数正负怎样控制左右两侧的升降；此处需要它判断函数经过候选点前后是否改变运动方向。使用时把候选点左右导数符号并排：\( +\to-\) 为极大，\(-\to+\) 为极小；最小例 \(x^2\) 在 \(0\) 左减右增，故 \(0\) 为极小值点。

\textbf{严格核心与条件：}若 \(x_0\) 是定义域内点，且存在邻域使 \(f(x)\le f(x_0)\)，则 \(x_0\) 为局部极大值点。Fermat 定理说：若 \(f\) 在内点 \(x_0\) 可导且在那里取极值，则 \(f'(x_0)=0\)。它只是必要条件；\(x^3\) 在 \(0\) 的导数为零却无极值。最稳的充分判据是导数变号。若 \(f'(x_0)=0\) 且 \(f''(x_0)>0\)，则为极小；\(f''(x_0)<0\) 则为极大；\(f''(x_0)=0\) 时判据失效。

\textbf{方法选择：}绝对值、分段或根式优先用一阶变号法；二阶导在驻点处非零时最快；参数题先找候选点，再按参数导致的“合并、消失、变号”分类。

\textbf{例 1（标准一阶判据）}\quad
\textbf{信号：}求多项式极值。\textbf{分析：}复用 K070 的符号表。
\textbf{完整解答：} \(f=x^3-3x\) 的 \(f'=3(x-1)(x+1)\)，符号为 \(+,-,+\)。故 \(x=-1\) 为极大值点，极大值 \(f(-1)=2\)；\(x=1\) 为极小值点，极小值 \(f(1)=-2\)。
\textbf{条件：}候选点均为可导内点。\textbf{易错：}把“点 \(-1\)”写成“极大值 \(-1\)”。\textbf{核验：}左右单调方向确实先增后减、先减后增。

\textbf{例 2（不可导极值）}\quad
\textbf{信号：}\(|x|\) 的尖点。\textbf{分析：}Fermat 不适用，但变号法可用。
\textbf{完整解答：} \(f=|x|\) 在 \(0\) 左侧导数为 \(-1\)，右侧为 \(1\)，故 \(0\) 是极小值点，极小值为 \(0\)。
\textbf{条件：}函数在候选点有定义；不要求该点可导。\textbf{易错：}由 \(f'(0)\) 不存在便排除极值。\textbf{核验：}\(|x|\ge0=f(0)\)。

\textbf{例 3（二阶判据退化）}\quad
\textbf{信号：}\(f'(0)=f''(0)=0\)。\textbf{分析：}二阶判据无结论，回到变号或定义。
\textbf{完整解答：}对 \(f=x^4\)，\(f'=4x^3\) 左负右正，所以 \(0\) 为极小值点；虽然 \(f''(0)=0\)，仍不影响结论。
\textbf{条件：}用的是一阶导数在去心邻域内的符号。\textbf{易错：}把 \(f''=0\) 写成“不是极值”。\textbf{核验：} \(x^4\ge0\)。

\textbf{例 4（参数退化）}\quad
\textbf{信号：} \(f=x^3-3ax\) 的极值随 \(a\) 变化。\textbf{分析：} \(f'=3(x^2-a)\) 的实根个数由 \(a\) 决定。
\textbf{完整解答：}若 \(a>0\)，在 \(x=-\sqrt a\) 处 \(+\to-\)，为极大值点；在 \(x=\sqrt a\) 处 \(-\to+\)，为极小值点。若 \(a=0\)，\(f'=3x^2\ge0\)，\(0\) 无极值；若 \(a<0\)，\(f'>0\)，无极值。
\textbf{条件：}参数为实数。\textbf{易错：}在 \(a=0\) 仍套出两个点。\textbf{核验：}边界 \(a=0\) 退化为 \(x^3\)。

\textbf{失分防线：}“驻点必为极值点”被 \(x^3\) 反驳；“极值点必满足 \(f'=0\)”被 \(|x|\) 反驳。快速检查候选点身份：是否内点、是否可导、导数是否真正变号。

\textbf{考场表达：}“在 \(x_0\) 左右，\(f'\) 由正变负（或负变正），故 \(x_0\) 为极大（或极小）值点，极值为 \(f(x_0)\)。”使用 Fermat 时必须先写“可导内点”。

\textbf{三级掌握标准：}基础过关：区分驻点与极值点；\(130+\)：会处理不可导点和二阶退化；满分能力：能完成参数分类并解释每个临界参数的几何意义。

\textbf{知识连接：}前置为 \knowledgeRef[MATH1-KN-CALC-0117]{K070}；局部候选将进入 \knowledgeRef[MATH1-KN-CALC-0119]{K072 全局最值}的候选表。
\end{knowledgeBox}

\studySubsection{calc-05-k072}{K072 闭区间最值：候选点完备清单}

\knowledgeAnchor[MATH1-KN-CALC-0119]{闭区间最大值最小值}
\noindent{\footnotesize\textbf{稳定引用：} \knowledgeRef[MATH1-KN-CALC-0119]{闭区间最大值最小值}}\par

\begin{knowledgeBox}
\textbf{定位与路线：}目标是求函数在指定集合上的全局最大值、最小值。核心不是只找极值，而是先保证“取得”，再把所有可能取得点列全。

\textbf{直接前置四要素桥：}\knowledgeRef[MATH1-KN-CALC-0118]{K071 局部极值}给出内部候选点；这里需要把局部候选与端点一起比较。使用规则是：先确认 \(f\in C[a,b]\)，再列端点、内部驻点、内部有定义但不可导点；最小例 \(x^2\) 在 \([-1,2]\) 的候选为 \(-1,0,2\)，比较值 \(1,0,4\)。

\textbf{严格核心与推导：}闭区间连续函数一定能取得最大值和最小值。若最值在内部可导点取得，Fermat 定理使其成为驻点；若内部不可导，也应列入；若不在内部，只能在端点。因此候选表完备。开区间、非闭集合或不连续函数不受该定理保证，此时必须用极限、单调性或定义直接判断“是否取到”。

\textbf{方法选择：}题目明确“在 \([a,b]\) 上”就用候选值比较法；绝对值/分段式补不可导点；实际最优化先把几何量写成单变量函数并确定闭区间；参数题先按候选点是否落在区间内分类。

\textbf{例 1（标准比较）}\quad
\textbf{信号：}多项式在闭区间上的最值。\textbf{分析：}端点加内部驻点。
\textbf{完整解答：} \(f=x^3-3x\) 在 \([-1,2]\) 上，内部驻点只有 \(x=1\)。比较 \(f(-1)=2,f(1)=-2,f(2)=2\)，最大值为 \(2\)，在 \(x=-1,2\) 取得；最小值为 \(-2\)，在 \(x=1\) 取得。
\textbf{条件：}多项式在闭区间连续。\textbf{易错：}漏端点或只报一个最大值点。\textbf{核验：}K071 的单调分段与比较结果一致。

\textbf{例 2（不可导候选）}\quad
\textbf{信号：}\(|x-1|\) 在区间内有尖点。\textbf{分析：}把 \(x=1\) 加入端点表。
\textbf{完整解答：}在 \([0,3]\) 上比较 \(f(0)=1,f(1)=0,f(3)=2\)，故最小值 \(0\) 在 \(1\) 取得，最大值 \(2\) 在 \(3\) 取得。
\textbf{条件：}函数连续，虽在 \(1\) 不可导。\textbf{易错：}只解 \(f'=0\) 得到空集。\textbf{核验：}距离函数的最小距离显然为零。

\textbf{例 3（定义域与候选过滤）}\quad
\textbf{信号：} \(f=x+1/x\) 在 \([1/2,3]\) 上求最值。\textbf{分析：}驻点 \(x=1\) 落在区间内。
\textbf{完整解答：} \(f'=1-1/x^2\)，候选为 \(1/2,1,3\)。对应值 \(5/2,2,10/3\)，故最小值 \(2\) 在 \(1\) 取得，最大值 \(10/3\) 在 \(3\) 取得。
\textbf{条件：}区间避开 \(0\)，函数连续。\textbf{易错：}把 \(x=-1\) 也列入。\textbf{核验：}正数域上 \(x+1/x\ge2\)。

\textbf{例 4（参数候选进入区间）}\quad
\textbf{信号：} \(f=x^2-2ax\) 在 \([0,1]\) 上求最小值。\textbf{分析：}驻点 \(x=a\) 是否在区间由参数决定。
\textbf{完整解答：}若 \(a\le0\)，函数不减，最小值 \(0\)；若 \(0<a<1\)，最小值 \(f(a)=-a^2\)；若 \(a\ge1\)，函数不增，最小值 \(f(1)=1-2a\)。
\textbf{条件：}所有参数下函数均连续。\textbf{易错：}不检查 \(x=a\) 是否属于 \([0,1]\)。\textbf{核验：}配方 \(f=(x-a)^2-a^2\) 给出同一投影结论。

\textbf{失分防线：}只写“解 \(f'=0\)”会漏端点和不可导点；开区间上的 \(f(x)=x\) 没有最大值，说明上确界不等于最大值。快速检查四格：连续闭区间、左端点、右端点、内部两类点。

\textbf{考场表达：}“因 \(f\in C[a,b]\)，最值必能取得。内部候选由 \(f'=0\) 及不可导点给出。列表比较……，所以……。”

\textbf{三级掌握标准：}基础过关：候选表无遗漏；\(130+\)：会处理参数、分段与实际建模；满分能力：能在非标准集合上先判断取得性，再给最大/最小或仅给上/下确界。

\textbf{知识连接：}候选筛选依赖 \knowledgeRef[MATH1-KN-CALC-0118]{K071}；与 \knowledgeRef[MATH1-KN-CALC-0138]{K091 积分估计}中的 \(m\le f\le M\) 直接连接。
\end{knowledgeBox}

\studySubsection{calc-05-k073}{K073 凹凸性与拐点：二阶导数判号}

\knowledgeAnchor[MATH1-KN-CALC-0120]{凹凸性与拐点}
\noindent{\footnotesize\textbf{稳定引用：} \knowledgeRef[MATH1-KN-CALC-0120]{凹凸性与拐点}}\par

\begin{knowledgeBox}
\textbf{定位与路线：}一阶导数描述升降，二阶导数描述斜率怎样变化。本讲统一用“\(f''\) 符号＋是否变号”表述，避免不同教材对“凹/凸”中文名称的相反约定。

\textbf{直接前置四要素桥：}\knowledgeRef[MATH1-KN-CALC-0117]{K070 单调区间判定}把导数符号变成函数的增减；此处对 \(f'\) 再用一次同一思想。二阶导数 \(f''=(f')'\) 是斜率的变化率，本题需要它判断曲线向哪侧弯。使用时找 \(f''=0\)、\(f''\) 不存在及定义域断点并判号；最小例 \(x^2\) 有 \(f''=2>0\)，斜率持续增加，图形按本书约定为凹。

\textbf{严格核心与条件：}在区间内若 \(f''>0\)，则 \(f'\) 严格增加，图形称凹；若 \(f''<0\)，图形称凸。拐点是曲线上凹凸性发生改变的点，必须写坐标 \((x_0,f(x_0))\)。方程 \(f''(x_0)=0\) 只提供候选；\(x^4\) 在 \(0\) 两侧均有 \(f''>0\)，故不是拐点。

\textbf{方法选择：}能求二阶导时直接判号；二阶导不存在时可比较 \(f'\) 的单调性或用弦线定义。参数题关注二阶导根的重合和越界。

\textbf{例 1（标准变号）}\quad
\textbf{信号：} \(f=x^3\)。\textbf{分析：} \(f''=6x\) 在 \(0\) 变号。
\textbf{完整解答：}在 \((-\infty,0)\) 上 \(f''<0\)，图形凸；在 \((0,\infty)\) 上 \(f''>0\)，图形凹；\((0,0)\) 为拐点。
\textbf{条件：}函数处处二阶可导。\textbf{易错：}只报横坐标 \(0\)。\textbf{核验：}斜率 \(3x^2\) 先减后增。

\textbf{例 2（零点不是拐点）}\quad
\textbf{信号：} \(f=x^4\) 且 \(f''(0)=0\)。\textbf{分析：}必须查两侧符号。
\textbf{完整解答：} \(f''=12x^2\ge0\)，在 \(0\) 两侧均为正，故全图形凹，\((0,0)\) 不是拐点。
\textbf{条件：}二阶导存在。\textbf{易错：}解 \(f''=0\) 后直接定拐点。\textbf{核验：}任一切线都位于图形下方。

\textbf{例 3（定义域先行）}\quad
\textbf{信号：} \(f=\ln x\)。\textbf{分析：}只在 \(x>0\) 讨论。
\textbf{完整解答：} \(f''=-1/x^2<0\)，故在 \((0,\infty)\) 上图形凸，无拐点。
\textbf{条件：}定义域为正半轴。\textbf{易错：}把 \(x=0\) 当成拐点。\textbf{核验：} \(x=0\) 不在曲线上，不能成为拐点。

\textbf{例 4（两个拐点）}\quad
\textbf{信号：} \(f=x^4-4x^3\)。\textbf{分析：} \(f''=12x(x-2)\) 易判号。
\textbf{完整解答：} \(f''>0\) 于 \((-\infty,0)\cup(2,\infty)\)，\(f''<0\) 于 \((0,2)\)，故拐点为 \((0,0)\) 与 \((2,-16)\)。
\textbf{条件：}多项式处处二阶可导。\textbf{易错：}漏算函数值。\textbf{核验：}符号顺序为 \(+,-,+\)，两点均真正变号。

\textbf{失分防线：}错误链“\(f''(x_0)=0\Rightarrow x_0\) 为拐点”被 \(x^4\) 反驳；定义域断点不在曲线上，不能叫拐点。快速检查：“在曲线上吗？两侧凹凸变了吗？”

\textbf{考场表达：}“\(f''=\cdots\)，其符号为……，故图形在……上 \(f''>0\)（凹），在……上 \(f''<0\)（凸）；因在 \(x_0\) 两侧变号，拐点为 \((x_0,f(x_0))\)。”

\textbf{三级掌握标准：}基础过关：二阶导符号表和拐点坐标正确；\(130+\)：能处理不可二阶可导候选；满分能力：能把凹凸性用于切线/弦线不等式与参数分类。

\textbf{知识连接：}与 \knowledgeRef[MATH1-KN-CALC-0117]{K070 单调性}同构，只是把对象从 \(f\) 换成 \(f'\)；是 \knowledgeRef[MATH1-KN-CALC-0122]{K075 作图}的关键层。
\end{knowledgeBox}

\studySubsection{calc-05-k074}{K074 渐近线：水平、垂直与斜线的完整扫描}

\knowledgeAnchor[MATH1-KN-CALC-0121]{渐近线}
\noindent{\footnotesize\textbf{稳定引用：} \knowledgeRef[MATH1-KN-CALC-0121]{渐近线}}\par

\begin{knowledgeBox}
\textbf{定位与路线：}渐近线描述曲线靠近定义域边界或趋向无穷时的线性骨架。扫描顺序为：所有有限定义域断点查垂直线；\(x\to+\infty\) 与 \(x\to-\infty\) 分别查水平线和斜线。

\textbf{直接前置四要素桥：}\knowledgeRef[MATH1-KN-CALC-0075]{K026 无穷处与无穷函数极限}描述变量逼近而不要求到达；本题用它衡量曲线到直线的距离是否趋零。使用时水平线查 \(f(x)\to b\)，垂直线查一侧 \(|f(x)|\to\infty\)，斜线查 \(f(x)-(ax+b)\to0\)；最小例 \(1/x\) 在 \(x\to0^\pm\) 无界，故 \(x=0\) 是垂直渐近线。

\textbf{严格核心与推导：}斜渐近线 \(y=ax+b\) 的必要计算为
\[
a=\lim_{x\to\pm\infty}\frac{f(x)}x,\qquad
b=\lim_{x\to\pm\infty}\bigl(f(x)-ax\bigr),
\]
两极限均须有限，通常 \(a\ne0\)。正负无穷可能得到不同直线；垂直线只需某个单侧极限绝对值无穷。

\textbf{方法选择：}有理函数先比较次数；根式差先有理化；含反三角函数可用已知极限或等价展开。算出水平线后无需再把同一端写成斜率 \(a=0\) 的斜线。

\textbf{例 1（有理函数）}\quad
\textbf{信号：}分母在 \(x=1\) 为零，分子不为零。\textbf{分析：}查垂直线并作长除。
\textbf{完整解答：} \(f=(2x+1)/(x-1)=2+3/(x-1)\)。故 \(x=1\) 为垂直渐近线，\(y=2\) 为 \(x\to\pm\infty\) 的水平渐近线。
\textbf{条件：} \(x\ne1\)。\textbf{易错：}只写水平线。\textbf{核验：} \(f-2=3/(x-1)\to0\)。

\textbf{例 2（两端不同）}\quad
\textbf{信号：} \(f=\sqrt{x^2+1}-x\)。\textbf{分析：}正端有理化，负端求斜线。
\textbf{完整解答：}当 \(x\to+\infty\)，\(f=1/(\sqrt{x^2+1}+x)\to0\)，故 \(y=0\)。当 \(x\to-\infty\)，
\[
a=\lim f/x=-2,\qquad b=\lim(f+2x)=\lim(\sqrt{x^2+1}+x)=0,
\]
故斜渐近线为 \(y=-2x\)。
\textbf{条件：}函数处处有定义。\textbf{易错：}用 \(\sqrt{x^2}=x\) 处理负无穷。\textbf{核验：}\(\sqrt{x^2}=|x|\)。

\textbf{例 3（单侧垂直线）}\quad
\textbf{信号：} \(f=\ln x\)，定义域边界为 \(0\)。\textbf{分析：}查 \(x\to0^+\)。
\textbf{完整解答：}\(\ln x\to-\infty\)，所以 \(x=0\) 是垂直渐近线；在 \(+\infty\) 处既无水平线，且 \(a=\lim\ln x/x=0\)、\(b=\lim\ln x=\infty\)，故无斜线。
\textbf{条件：}只有右侧可逼近 \(0\)。\textbf{易错：}要求两侧都无穷才承认垂直线。\textbf{核验：}定义域本就没有 \(0\) 左侧。

\textbf{例 4（两条斜渐近线）}\quad
\textbf{信号：} \(f=x\arctan x\)。\textbf{分析：}正负无穷分算。
\textbf{完整解答：}由 \(\arctan x=\pi/2-1/x+o(1/x)\)（\(x\to+\infty\)），得 \(f=(\pi/2)x-1+o(1)\)；负端有 \(\arctan x=-\pi/2-1/x+o(1/x)\)，得 \(f=-(\pi/2)x-1+o(1)\)。故两条斜线分别为 \(y=\frac\pi2x-1\)、\(y=-\frac\pi2x-1\)。
\textbf{条件：}展开只用于对应端。\textbf{易错：}默认两端同线。\textbf{核验：} \(f\) 为偶函数，两条线关于 \(y\) 轴对称。

\textbf{失分防线：}只算 \(+\infty\) 会漏线；垂直线只看分母为零会把可约因子误判，例如 \((x^2-1)/(x-1)\) 在 \(1\) 可去间断但不无界。快速检查以极限定义为准。

\textbf{考场表达：}分别写清逼近方向；斜线必须展示 \(a,b\) 两个极限，最后核验 \(\lim[f-(ax+b)]=0\)。

\textbf{三级掌握标准：}基础过关：三类公式无混淆；\(130+\)：两端、单侧与可去因子均不漏；满分能力：能在根式和反三角结构中稳定提取线性主部。

\textbf{知识连接：}渐近线与 \knowledgeRef[MATH1-KN-CALC-0122]{K075 图形描绘}直接相连，并复用函数极限与定义域。
\end{knowledgeBox}

\studySubsection{calc-05-k075}{K075 函数图形描绘：把所有判据汇成一张草图}

\knowledgeAnchor[MATH1-KN-CALC-0122]{函数图形描绘}
\noindent{\footnotesize\textbf{稳定引用：} \knowledgeRef[MATH1-KN-CALC-0122]{函数图形描绘}}\par

\begin{knowledgeBox}
\textbf{定位与路线：}作图不是凭点连线，而是把定义域、对称性、截距、符号、单调极值、凹凸拐点、渐近线与端点趋势按顺序汇总。

\textbf{直接前置四要素桥：}\knowledgeRef[MATH1-KN-CALC-0117]{K070} 给升降，\knowledgeRef[MATH1-KN-CALC-0120]{K073} 给弯曲，\knowledgeRef[MATH1-KN-CALC-0121]{K074} 给远端骨架；本讲需要三者共同约束图形。使用时先画定义域隔断与渐近线，再放关键点和箭头；最小例 \(1/x\) 的定义域断点 \(0\) 与两条渐近线已经决定双曲线不能跨轴连接。

\textbf{严格核心与方法：}推荐清单：
\[
D_f\to\text{奇偶/周期}\to\text{零点与符号}\to f'\to f''\to\text{渐近线}\to\text{关键点}.
\]
草图只需拓扑结构正确，不追求比例。选图题可反向用零点、端行为、极值个数快速排除。

\textbf{例 1（多项式 S 形）}\quad
\textbf{信号：} \(f=x^3-3x\)。\textbf{分析：}奇函数，已有一阶导结论。
\textbf{完整解答：}零点为 \(-\sqrt3,0,\sqrt3\)；在 \(-1\) 极大值 \(2\)，在 \(1\) 极小值 \(-2\)；\(f''=6x\)，原点为拐点；无渐近线。按奇对称连接即可。
\textbf{条件：}定义域全实数。\textbf{易错：}极值点与零点次序画反。\textbf{核验：}所有结构关于原点中心对称。

\textbf{例 2（有水平渐近线）}\quad
\textbf{信号：} \(f=x/(1+x^2)\)。\textbf{分析：}奇函数，远端趋零。
\textbf{完整解答：} \(f'=(1-x^2)/(1+x^2)^2\)，故在 \([-1,1]\) 增、两侧减，极小点 \((-1,-1/2)\)、极大点 \((1,1/2)\)；\(y=0\) 是两端水平渐近线，且过原点。
\textbf{条件：}分母恒正。\textbf{易错：}曲线跨过渐近线后错误封闭。\textbf{核验：}奇对称且函数符号与 \(x\) 相同。

\textbf{例 3（定义域边界与极值）}\quad
\textbf{信号：} \(f=\ln x/x\)。\textbf{分析：}先定 \(x>0\)，再查两端。
\textbf{完整解答：} \(x\to0^+\) 时 \(f\to-\infty\)，故 \(x=0\) 为垂直渐近线；\(x\to\infty\) 时 \(f\to0^+\)，故 \(y=0\) 为水平渐近线。\(f'=(1-\ln x)/x^2\)，在 \(x=e\) 取极大值 \(1/e\)，零点为 \(x=1\)。
\textbf{条件：}只画右半平面。\textbf{易错：}把 \(y=0\) 当作全程不能穿越；曲线在 \(x=1\) 穿过它。\textbf{核验：}导数在 \(e\) 左正右负。

\textbf{例 4（指数衰减骨架）}\quad
\textbf{信号：} \(f=(x+1)e^{-x}\)。\textbf{分析：}指数恒正，零点和符号由 \(x+1\) 决定。
\textbf{完整解答：} \(f'= -xe^{-x}\)，故在 \((-\infty,0]\) 增、\([0,\infty)\) 减，极大值 \(1\)；零点 \(-1\)；\(x\to+\infty\) 时 \(f\to0^+\)，而 \(x\to-\infty\) 时 \(f\to-\infty\)。\(f''=(x-1)e^{-x}\)，拐点为 \((1,2/e)\)。
\textbf{条件：}定义域全实数。\textbf{易错：}误以为负无穷端也趋零。\textbf{核验：}符号、单调与端行为能够无冲突连接。

\textbf{失分防线：}跨定义域断点连线、把渐近线画成绝不可穿越、漏写左右端趋势都属结构错误。快速检查：每一段定义域必须有自己的左端、右端箭头。

\textbf{考场表达：}先给性质表，再画草图；标出坐标、渐近线和箭头。若用于根的个数，只需保留决定交点数的结构。

\textbf{三级掌握标准：}基础过关：按清单完成无冲突草图；\(130+\)：能从图像反推根数和面积分段；满分能力：参数改变时能定位拓扑变化的临界参数。

\textbf{知识连接：}综合 \knowledgeRef[MATH1-KN-CALC-0117]{K070}、\knowledgeRef[MATH1-KN-CALC-0118]{K071}、\knowledgeRef[MATH1-KN-CALC-0120]{K073}、\knowledgeRef[MATH1-KN-CALC-0121]{K074}，并服务于定积分面积分区。
\end{knowledgeBox}

\studySubsection{calc-05-k076}{K076 弧微分与曲率：显函数和参数曲线}

\knowledgeAnchor[MATH1-KN-CALC-0123]{弧微分与曲率公式}
\noindent{\footnotesize\textbf{稳定引用：} \knowledgeRef[MATH1-KN-CALC-0123]{弧微分与曲率公式}}\par

\begin{knowledgeBox}
\textbf{定位与路线：}曲率 \(\kappa\) 衡量单位弧长内切线方向转动得多快；直线曲率为 \(0\)，半径越小的圆曲率越大。

\textbf{直接前置四要素桥：}\knowledgeRef[MATH1-KN-CALC-0120]{K073 二阶导数判号}已经使用一、二阶导数描述图形形态；此处把一、二阶导数扩展到参数曲线的速度与方向变化。弧微分是曲线上微小位移的长度。对参数曲线 \((x(t),y(t))\)，速度向量为 \((x',y')\)，故
\[
ds=\sqrt{x'^2+y'^2}\,dt.
\]
本题需要用它把“切线角对参数的变化”换成“对弧长的变化”；使用前检查速度向量非零。最小例直线 \((t,0)\) 有 \(ds=dt\)，方向不转，曲率为零。

\textbf{严格核心与推导：}定义 \(\kappa=|d\theta/ds|\)。对显函数 \(y=f(x)\)，由 \(\tan\theta=y'\) 得
\[
\kappa=\frac{|y''|}{[1+(y')^2]^{3/2}}.
\]
对正则参数曲线（\(x'^2+y'^2\ne0\)），
\[
\kappa=\frac{|x'y''-y'x''|}{(x'^2+y'^2)^{3/2}}.
\]
绝对值保证曲率是大小；分母的 \(3/2\) 次方来自弧长归一化。

\textbf{方法选择：}显函数直接用第一式；参数形式不要先强行消参；若速度为零，公式不可硬套，应换参数或回到几何。

\textbf{例 1（抛物线顶点）}\quad
\textbf{信号：} \(y=x^2\) 在 \(0\) 的曲率。\textbf{分析：}显函数公式。
\textbf{完整解答：} \(y'=2x,y''=2\)，故 \(\kappa(0)=2/[1+0]^{3/2}=2\)。
\textbf{条件：}二阶导存在。\textbf{易错：}把分母写成一次方。\textbf{核验：}顶点附近弯曲最明显。

\textbf{例 2（圆的校准）}\quad
\textbf{信号：}上半圆 \(y=\sqrt{R^2-x^2}\) 在 \(0\)。\textbf{分析：}代入显函数式。
\textbf{完整解答：} \(y'(0)=0,y''(0)=-1/R\)，故 \(\kappa(0)=1/R\)。
\textbf{条件：} \(R>0\)。\textbf{易错：}漏绝对值得到负曲率。\textbf{核验：}圆上曲率应恒为 \(1/R\)。

\textbf{例 3（参数圆）}\quad
\textbf{信号：} \(x=a\cos t,y=a\sin t\)。\textbf{分析：}参数公式最短。
\textbf{完整解答：}分子 \(|x'y''-y'x''|=a^2\)，分母 \((a^2)^{3/2}=a^3\)，故 \(\kappa=1/a\)。
\textbf{条件：} \(a>0\)，速度大小为 \(a\ne0\)。\textbf{易错：}把二阶参数导数相除。\textbf{核验：}与圆的几何结论一致。

\textbf{例 4（参数抛物线）}\quad
\textbf{信号：} \(x=t,y=t^2\) 在 \(t=1\)。\textbf{分析：} \(x'=1\) 使正则性显然。
\textbf{完整解答：} \(x'=1,x''=0,y'=2t,y''=2\)，故
\[
\kappa(1)=\frac2{(1+4)^{3/2}}=\frac2{5^{3/2}}.
\]
\textbf{条件：}速度从不为零。\textbf{易错：}漏掉参数导数。\textbf{核验：}消参得 \(y=x^2\)，代 \(x=1\) 同值。

\textbf{失分防线：}三大错误是漏绝对值、分母指数错、速度为零仍套式。快速量纲检查：圆半径放大 \(c\) 倍，曲率应缩小为 \(1/c\)。

\textbf{考场表达：}先列一二阶导与正则性，再代公式；参数式保留“分子行列式、分母速度三次方”的结构。

\textbf{三级掌握标准：}基础过关：两套公式准确；\(130+\)：会检查正则性并用消参交叉核验；满分能力：能从 \(\kappa=|d\theta/ds|\) 解释公式和尺度变化。

\textbf{知识连接：}直接后续是 \knowledgeRef[MATH1-KN-CALC-0124]{K077 曲率半径与曲率圆}；弧微分也连接积分学中的弧长。
\end{knowledgeBox}

\studySubsection{calc-05-k077}{K077 曲率圆与曲率半径：从弯曲程度到密切圆}

\knowledgeAnchor[MATH1-KN-CALC-0124]{曲率圆、曲率半径}
\noindent{\footnotesize\textbf{稳定引用：} \knowledgeRef[MATH1-KN-CALC-0124]{曲率圆、曲率半径}}\par

\begin{knowledgeBox}
\textbf{定位与路线：}曲率圆是在给定点与曲线具有二阶接触的圆；其半径称曲率半径。它把抽象曲率变成“此处最贴合的圆有多大”。

\textbf{直接前置四要素桥：}\knowledgeRef[MATH1-KN-CALC-0123]{K076 曲率}给出单位弧长的转角变化率；这里需要取倒数得到长度尺度。使用规则为 \(\rho=1/\kappa\)（仅 \(\kappa>0\)）；最小例半径 \(R\) 的圆有 \(\kappa=1/R\)，所以 \(\rho=R\)。

\textbf{严格核心与条件：}在 \(\kappa\ne0\) 的正则点，曲率圆半径
\[
\rho=\frac1\kappa,
\]
圆心沿主法线方向距该点 \(\rho\)。对图形 \(y=f(x)\) 且 \(y''\ne0\)，带符号的圆心可写为
\[
\left(x-\frac{y'(1+y'^2)}{y''},\qquad
y+\frac{1+y'^2}{y''}\right).
\]
若 \(\kappa=0\)，曲率半径可理解为无穷大，但不存在普通有限曲率圆。

\textbf{方法选择：}只问半径先算 \(\kappa\) 再取倒数；问圆方程才求中心。中心公式需保留 \(y''\) 符号，不能把曲率绝对值直接塞入方向。

\textbf{例 1（抛物线顶点）}\quad
\textbf{信号：} \(y=x^2\) 在原点的曲率圆。\textbf{分析：}K076 得 \(\kappa=2\)。
\textbf{完整解答：} \(\rho=1/2\)。中心公式给 \(C=(0,1/2)\)，故曲率圆为 \(x^2+(y-1/2)^2=1/4\)。
\textbf{条件：} \(y''(0)=2\ne0\)。\textbf{易错：}把中心放到下方。\textbf{核验：}圆在原点的切线也是 \(y=0\)。

\textbf{例 2（圆自身）}\quad
\textbf{信号：}半径 \(R\) 的圆。\textbf{分析：}曲率恒定。
\textbf{完整解答：}\(\kappa=1/R,\rho=R\)，曲率圆就是原圆。
\textbf{条件：} \(R>0\)。\textbf{易错：}写成 \(\rho=\kappa\)。\textbf{核验：}二阶接触在圆的每一点成立。

\textbf{例 3（零曲率边界）}\quad
\textbf{信号：} \(y=x^3\) 在原点。\textbf{分析：} \(y''(0)=0\)。
\textbf{完整解答：}\(\kappa(0)=0\)，故无有限曲率半径和普通曲率圆；不能代中心公式，因为分母 \(y''\) 为零。
\textbf{条件：}曲线正则，但曲率为零。\textbf{易错：}写出“半径为 \(0\)”。\textbf{核验：}原点处曲线二阶近似为切线。

\textbf{例 4（椭圆端点）}\quad
\textbf{信号：} \(x=a\cos t,y=b\sin t\) 在 \(t=0\)。\textbf{分析：}参数曲率式。
\textbf{完整解答：}分子为 \(ab\)，速度为 \(b\)，故 \(\kappa=a/b^2\)，\(\rho=b^2/a\)。
\textbf{条件：} \(a,b>0\)。\textbf{易错：}把端点速度误写为 \(a\)。\textbf{核验：}当 \(a=b=R\) 退化为 \(\rho=R\)。

\textbf{失分防线：}\(\kappa=0\) 时取倒数不是零而是无有限值；中心方向由有符号法线决定。快速检查：算得的半径必须为正，并在圆的特例上回归 \(R\)。

\textbf{考场表达：}“先由曲率公式得 \(\kappa=\cdots>0\)，故 \(\rho=1/\kappa\)；若求曲率圆，再由主法线方向得圆心……。”

\textbf{三级掌握标准：}基础过关：曲率与半径不混；\(130+\)：会算简单曲率圆；满分能力：能处理 \(\kappa=0\) 的退化并用圆/椭圆交叉核验。

\textbf{知识连接：}前置为 \knowledgeRef[MATH1-KN-CALC-0123]{K076}；下一节 \knowledgeRef[MATH1-KN-CALC-0125]{K078}只作边界辨识，不把渐屈线当统考主线。
\end{knowledgeBox}

\studySubsection{calc-05-k078}{K078 大纲边界：曲率中心、渐屈线与方程近似解}

\knowledgeAnchor[MATH1-KN-CALC-0125]{曲率中心、渐屈线与方程近似解}
\noindent{\footnotesize\textbf{稳定引用：} \knowledgeRef[MATH1-KN-CALC-0125]{曲率中心、渐屈线与方程近似解}}\par

\begin{knowledgeBox}
\textbf{定位与路线：}本节点的主要任务是边界识别：曲率、曲率圆和曲率半径属于本章应掌握内容；渐屈线、渐伸线及二分法、Newton 迭代不作为统考数学一专项训练。下面只给理解性桥接，不能用它们替代单调性、零点定理或误差论证。

\textbf{直接前置四要素桥：}\knowledgeRef[MATH1-KN-CALC-0124]{K077 曲率圆}中的圆心是曲线上每一点对应的几何对象；若让点移动，圆心轨迹称渐屈线，这说明该术语“是什么、为何出现、怎样构造”。最小例是圆：每点的曲率中心都是原圆心，所以其渐屈线退化为一点。方程近似解则以切线代替曲线，如 Newton 步 \(x_{n+1}=x_n-f(x_n)/f'(x_n)\)。

\textbf{严格核心与边界：}“会辨认”不等于“考场可无条件调用”。Newton 公式要求当前点处 \(f'\ne0\)，其收敛还需初值和函数形状条件；二分法要求连续且端点异号。统考存在唯一性仍应优先写连续性、零点定理和严格单调性。

\textbf{方法选择：}若题面只问曲率圆，停在 K077；若学校自命题明确给出迭代式，按题设计算；若只问根的个数，不引入数值迭代。

\textbf{例 1（范围辨识）}\quad
\textbf{信号：}在“曲率、曲率半径、渐屈线、Newton 迭代”中选择统考主线。
\textbf{分析：}按本节点边界分类。\textbf{完整解答：}应掌握前两项；后两项只作教材扩展辨识，不安排专项公式背诵。
\textbf{条件：}结论限于本仓库的数学一备考口径。\textbf{易错：}把教材出现等同于必考。\textbf{核验：}主线知识已由 K076、K077 闭合。

\textbf{例 2（圆的渐屈线直觉）}\quad
\textbf{信号：}问半径 \(R\) 圆的曲率中心轨迹。\textbf{分析：}每点曲率圆都是原圆。
\textbf{完整解答：}每点曲率中心均为圆心 \(O\)，故轨迹退化为单点 \(O\)。
\textbf{条件：} \(R>0\)。\textbf{易错：}误答为原圆。\textbf{核验：}曲率半径始终等于 \(R\) 且主法线都指向 \(O\)。

\textbf{例 3（Newton 一步仅作理解）}\quad
\textbf{信号：}用切线近似 \(\sqrt2\)。\textbf{分析：}对 \(f(x)=x^2-2\) 从 \(x_0=1\) 作一次切线。
\textbf{完整解答：} \(x_1=x_0-f(x_0)/f'(x_0)=1-(-1)/2=3/2\)。
\textbf{条件：} \(f'(1)=2\ne0\)；这只是一轮近似，不是精确根。\textbf{易错：}写成 \(\sqrt2=3/2\)。\textbf{核验：} \(1.4<\sqrt2<1.5\)，近似方向合理。

\textbf{例 4（证明与数值法分工）}\quad
\textbf{信号：}证明 \(x+\ln x=0\) 有唯一根并给粗区间。
\textbf{分析：}证明用连续单调，区间用异号，不需 Newton。
\textbf{完整解答：} \(g'=1+1/x>0\) 于 \(x>0\)，且 \(g(1/2)<0<g(1)\)，故在 \((1/2,1)\) 恰有一根。
\textbf{条件：} \(g\) 在该区间连续。\textbf{易错：}用若干迭代数值代替唯一性证明。\textbf{核验：}结论同时含存在、唯一与定位。

\textbf{失分防线：}未经条件检查直接写 Newton 迭代不是证明；近似值不能当精确等式。快速检查：题目要的是“证明、精确计算”还是“数值近似”。

\textbf{考场表达：}统考主线只使用题内已建立的定理；扩展术语若题设未定义，不用作关键跳步。

\textbf{三级掌握标准：}基础过关：能区分主线与扩展；\(130+\)：不会让数值近似替代严格证明；满分能力：题设明确给迭代时能核验导数非零和误差方向，但不额外扩展。

\textbf{知识连接：}几何部分回到 \knowledgeRef[MATH1-KN-CALC-0124]{K077}；方程根的严格判定回到 \knowledgeRef[MATH1-KN-CALC-0117]{K070}。
\end{knowledgeBox}
